%15 min preso!
\documentclass[xcolor=table,aspectratio=169]{beamer}
\usepackage{beamerthemesplit}
\usepackage{wrapfig}
\usetheme{SPbGU}
\usepackage{pdfpages}
\usepackage{amsmath}
\usepackage{cmap}
\usepackage[T2A]{fontenc}
\usepackage[utf8]{inputenc}
\usepackage[english]{babel}
\usepackage{indentfirst}
\usepackage{amsmath}
\usepackage{tikz}
\usepackage{multirow}
\usepackage[noend]{algpseudocode}
\usepackage{algorithm}
\usepackage{algorithmicx}
\usepackage{fancyvrb}
\usepackage{hyperref} 
\definecolor{links}{HTML}{2A1B81}
\hypersetup{colorlinks,linkcolor=,urlcolor=links}
\usetikzlibrary{calc}
\usetikzlibrary{shapes, backgrounds}
\usetikzlibrary{arrows,automata}
\usetikzlibrary{positioning}
\usetikzlibrary{fit}
\usetikzlibrary{shapes.callouts}
\usetikzlibrary{shapes.misc}
\usepackage{xparse}
\usepackage{fontawesome}

\usepackage{etoolbox,refcount}
\usepackage{multicol}

\usepackage{tabularx}
\newcolumntype{Y}{>{\raggedleft\arraybackslash}X}

\renewcommand{\thealgorithm}{}

\newtheorem{mytheorem}{Theorem}
\renewcommand{\thealgorithm}{}

\newcommand{\tikzmark}[1]{\tikz[overlay,remember picture] \node (#1) {};}
\def\Put(#1,#2)#3{\leavevmode\makebox(0,0){\put(#1,#2){#3}}}

\newcommand{\ltz}{$< 1$}

\tikzset{
    state/.style={
           rectangle,
           rounded corners,
           draw=black, very thick,
           minimum height=2em,
           inner sep=2pt,
           text centered,
           },
}

\tikzset{
    invisible/.style={opacity=0,text opacity=0},
    visible on/.style={alt=#1{}{invisible}},
    alt/.code args={<#1>#2#3}{%
      \alt<#1>{\pgfkeysalso{#2}}{\pgfkeysalso{#3}} % \pgfkeysalso doesn't change the path
    },
}

\tikzset{cross/.style={cross out, draw=black, minimum size=2*(#1-\pgflinewidth), inner sep=0pt, outer sep=0pt, ultra thick},
%default radius will be 1pt. 
cross/.default={1pt}}

\NewDocumentCommand{\mycallout}{r<> O{opacity=0.8,text opacity=1} m m m}{%
\tikz[remember picture, overlay]\node[align=center, fill=cyan!20, text width=#5cm,
#2,visible on=<#1>, rounded corners,
draw,rectangle callout,anchor=pointer,callout relative pointer={(290:0.5cm)}]
at (#3) {#4};
}

\NewDocumentCommand{\mycalloutR}{r<> O{opacity=0.8,text opacity=1} m m m}{%
\tikz[remember picture, overlay]\node[align=center, fill=cyan!20, text width=#5cm,
#2,visible on=<#1>, rounded corners,
draw,rectangle callout,anchor=pointer,callout relative pointer={(30:0.8cm)}]
at (#3) {#4};
}

\newcommand\colR{\cellcolor{red!20}}
\newcommand\colB{\cellcolor{blue!20}}
\newcommand\colG{\cellcolor{green!20}}
\definecolor{Gray}{gray}{0.8}

%callout relative pointer={(230:0.5cm)}]

\newcounter{countitems}
\newcounter{nextitemizecount}
\newcommand{\setupcountitems}{%
  \stepcounter{nextitemizecount}%
  \setcounter{countitems}{0}%
  \preto\item{\stepcounter{countitems}}%
}
\makeatletter
\newcommand{\computecountitems}{%
  \edef\@currentlabel{\number\c@countitems}%
  \label{countitems@\number\numexpr\value{nextitemizecount}-1\relax}%
}
\newcommand{\nextitemizecount}{%
  \getrefnumber{countitems@\number\c@nextitemizecount}%
}
\newcommand{\previtemizecount}{%
  \getrefnumber{countitems@\number\numexpr\value{nextitemizecount}-1\relax}%
}
\makeatother    
\newenvironment{AutoMultiColItemize}{%
\ifnumcomp{\nextitemizecount}{>}{3}{\begin{multicols}{2}}{}%
\setupcountitems\begin{itemize}}%
{\end{itemize}%
\unskip\computecountitems\ifnumcomp{\previtemizecount}{>}{3}{\end{multicols}}{}}


\beamertemplatenavigationsymbolsempty

\title[Анализ графов на RISC-V]{Анализ графов и линейная алгебра в экосистеме RISC-V}
\subtitle{Обновление статуса}
\institute[СПбГУ]{
Санкт-Петербургский Государственный Университет
}

\author[Семён Григорьев]{Семён Григорьев}

\date{26 июня 2026}

\begin{document}
{
\begin{frame}[fragile]
  \begin{table}
  \centering  
  \begin{tabularx}{\linewidth}{XcX}
    \hfill
    & 
    & \hfill \includegraphics[height=1.4cm]{pictures/SPbSU_Logo.pdf}
  \end{tabularx}
  \end{table}
  \titlepage
\end{frame}
}

\begin{frame}[fragile]
  \frametitle{Вышел Vortex 3.0}
  \begin{itemize}
    \item Release notes: \url{https://github.com/vortexgpgpu/vortex/releases/tag/v3.0}
    \begin{itemize} 
      \item Все компоненты: RTL, рантайм, инструменты разработчика, \ldots
      \item Вошли наши исправления и доработки
    \end{itemize}
    \item Ключевые (для нас) обновления
    \begin{itemize} 
      \item Атомики в RTL
      \item Обновление до LLVM 20 и POCL 7.0
      \item Kernel-entry calling convention (no callee-saved spills)
      \begin{itemize} 
        \item В наших issue про сброс регистров ответили, что <<странное>> поведение перестало воспроизводиться
      \end{itemize}
    \end{itemize}
    \item[\faGears] \textbf{Надо обновлять и всё перепроверять}
  \end{itemize}
\end{frame}

\begin{frame}[fragile]
  \frametitle{Spla на Vortex}
    \begin{itemize}  
    \item[\faGears] Должно всё работать
    \begin{itemize}
      \item Считаются эксперименты в симуляторах
      \item Пока на старой версии с нашими фиксами
    \end{itemize}
    \item[\faGears] План:
    \begin{itemize}
      \item Закончить эксперименты на старой версии Vortex
      \item Что-то самое интересное перепроверить на новой версии
      \item Запуск на ПЛИС с новым Vortex
    \end{itemize}
    \item[\faGears] Пишем статью
  \end{itemize}
\end{frame}

\begin{frame}[fragile]
  \frametitle{Запуск GPGPU на Milk-V}
  \begin{itemize}
    \item[\faCheck] Старый AMD (Radeon HD 6450)
    \begin{itemize}
      \item Запускается <<из коробки>>
      \item Работают прикладные алгоритмы на основе Spla
    \end{itemize} 
    \item[\faGears] Новый AMD (RX 6750)
    \begin{itemize} 
      \item Должно быть возможно
      \item Нужно конфигурировть и собирать ядро
    \end{itemize}
    \item[\faGears] Nvidia
    \begin{itemize} 
      \item Должно быть возможно (скорее всего, через \href{https://nouveau.freedesktop.org/}{nouveau})
      \item Нужно конфигурировть и собирать ядро
    \end{itemize}
    \item[\faGears] Intel ARC
    \begin{itemize} 
      \item Должно быть возможно
      \item Нужно конфигурировть и собирать ядро
    \end{itemize}
  \end{itemize}
\end{frame}

\end{document}
